\paragraph{One virtual operation on every staircase window.}

The construction below formalizes the forward virtual/physical-operation
correspondence in the proof sketch of the two-dimensional corollary
\cite[corollary after Lemma~5, lines~2297--2444]{Molnar2018NormalPEPS}.

\begin{definition}[Partial state across a region cut]%
    \label{def:peps_region_partial_state}
    \lean{TNLean.PEPS.regionPartialState}
    \leanok
    For a region $R$ and a physical configuration $\tau$ on its complement,
    define the partial state $\psi^\tau_{A,R}$ on the physical space of $R$
    by
    \begin{align}
        \psi^\tau_{A,R}(\sigma)
         & =c_A(\omega_{R,\sigma,\tau}).
        \notag
    \end{align}
\end{definition}

\begin{theorem}[Partial states of equal closed states]%
    \label{thm:peps_region_partial_state_same_state}
    \lean{TNLean.PEPS.regionPartialState_sameState}
    \leanok
    \uses{def:peps_region_partial_state}
    If $A$ and $B$ generate the same closed state, then, for every region
    $R$ and every complement configuration $\tau$,
    $\psi^\tau_{A,R}=\psi^\tau_{B,R}$.
\end{theorem}
\begin{proof}
    \leanok
    Evaluate both functions at a physical configuration $\sigma$ on $R$ and
    apply the equality of the two closed-state coefficients at
    $\omega_{R,\sigma,\tau}$.
\end{proof}

\begin{theorem}[Blocked-tensor expansion of a partial state]%
    \label{thm:peps_region_partial_state_blocked_expansion}
    \lean{TNLean.PEPS.regionInteriorBondProd_smul_regionPartialState}
    \leanok
    \uses{def:peps_region_partial_state,
        def:peps_regionComplementBoundaryConfig}
    Let $P_A(R)$ be the product of the dimensions of the bonds internal to
    $R$. For every complement configuration $\tau$,
    \begin{align}
        P_A(R)\psi^\tau_{A,R}
         & =\sum_{\mu\in\partial R}
        T_{A,V\setminus R}(\bar\mu)(\tau)T_{A,R}(\mu).
        \label{eq:peps_partial_state_blocked_expansion}
    \end{align}
\end{theorem}
\begin{proof}
    \leanok
    Split the closed contraction across $R$. If $\mu$ and $\nu$ are boundary
    labels on $R$, the complement-boundary identification satisfies
    \begin{align}
        \bar\mu=\bar\nu
         & \quad\Longleftrightarrow\quad \mu=\nu.
        \notag
    \end{align}
    Hence the double boundary sum collapses to its diagonal, which gives the
    sum in~(\ref{eq:peps_partial_state_blocked_expansion}); each bond internal
    to $R$ contributes its dimension to the factor $P_A(R)$.
\end{proof}

\begin{definition}[Physical operator of a region insert]%
    \label{def:peps_physical_operator_of_region_insert}
    \lean{TNLean.PEPS.physicalOpOfRegionInsert}
    \leanok
    \uses{def:peps_region_out_endpoint_block_inverse}
    Let $R$ be an injective region and let
    $C=\{C(\mu)\}_{\mu\in\partial R}$ be any family in its physical
    space, indexed by the virtual boundary configurations.  If $L_{A,R}$
    is a left inverse of the blocked tensor map, define
    \begin{align}
        O_C
         & =\left(c\longmapsto\sum_{\mu}c_\mu C(\mu)\right)L_{A,R}.
        \notag
    \end{align}
\end{definition}

\begin{theorem}[Open-boundary realization of a region insert]%
    \label{thm:peps_physical_operator_of_region_insert_realizes}
    \lean{TNLean.PEPS.physicalOpOfRegionInsert_realizes}
    \leanok
    \uses{def:peps_physical_operator_of_region_insert}
    For every virtual boundary configuration $\mu$,
    \begin{align}
        O_C\bigl(T_{A,R}(\mu)\bigr) & =C(\mu).
        \notag
    \end{align}
    Thus the insert is obtained from the genuine region block by one
    physical operator before the region is contracted with its complement.
\end{theorem}
\begin{proof}
    \leanok
    The left inverse sends $T_{A,R}(\mu)$ to the standard basis vector
    indexed by $\mu$, and the displayed linear combination sends that basis
    vector to $C(\mu)$.
\end{proof}

\begin{definition}[Region insert of a physical operator]%
    \label{def:peps_region_insert_of_physical_operator}
    \lean{TNLean.PEPS.regionInsertOfPhysicalOp}
    \leanok
    Let $O$ be a physical operator on the physical space of a region $R$.
    Its action on the genuine open-boundary blocks defines the insert
    \begin{align}
        C^O_{A,R}(\mu) & =O\bigl(T_{A,R}(\mu)\bigr).
        \notag
    \end{align}
\end{definition}

\begin{theorem}[Closed state of a physical operator]%
    \label{thm:peps_closed_state_of_physical_operator}
    \lean{TNLean.PEPS.deformedRegionState_regionInsertOfPhysicalOp}
    \leanok
    \uses{def:peps_region_insert_of_physical_operator,
        thm:peps_region_partial_state_blocked_expansion}
    For a complement physical configuration $\tau$, let
    $\psi^\tau_{A,R}$ be the partial state across the cut. Then
    \begin{align}
        \Psi_{A,R,C^O_{A,R}}(\sigma,\tau)
         & =\bigl[O\bigl(P_A(R)\psi^\tau_{A,R}\bigr)\bigr](\sigma).
        \label{eq:peps_physical_op_partial_state}
    \end{align}
\end{theorem}
\begin{proof}
    \leanok
    Apply $O$ to~(\ref{eq:peps_partial_state_blocked_expansion}) and evaluate
    at $\sigma$. This proves
    (\ref{eq:peps_physical_op_partial_state}).
\end{proof}

\begin{theorem}[Physical operator transferred between equal closed states]%
    \label{thm:peps_physical_operator_same_state_transfer}
    \lean{TNLean.PEPS.deformedRegionState_regionInsertOfPhysicalOp_sameState}
    \leanok
    \uses{thm:peps_closed_state_of_physical_operator,
        thm:peps_region_partial_state_same_state}
    If $A$ and $B$ generate the same closed state, then
    \begin{align}
        P_B(R)\Psi_{A,R,C^O_{A,R}}(\sigma,\tau)
         & =P_A(R)\Psi_{B,R,C^O_{B,R}}(\sigma,\tau).
        \label{eq:peps_physical_op_same_state}
    \end{align}
    The two virtual boundary spaces need not be identified.
\end{theorem}
\begin{proof}
    \leanok
    Equality of the closed states gives
    $\psi^\tau_{A,R}=\psi^\tau_{B,R}$. Apply
    (\ref{eq:peps_physical_op_partial_state}) to both tensors and
    cross-multiply the two interior-bond products.
\end{proof}

\begin{definition}[Interior-bond region insert]%
    \label{def:peps_interior_bond_inserted_region_insert}
    \lean{TNLean.PEPS.interiorBondInsertedRegionInsert}
    \leanok
    \uses{def:peps_bond_inserted_region_insert, def:peps_corner_extension}
    Let $e=\{u,w\}$ be an ordered edge whose two endpoints lie in a region
    $R$, with $u$ its ordered left endpoint, and let
    $X\in\MN{D_A(e)}$.  Write $U=\{u\}$ and let $I_{A,U,e,X}$ be the
    boundary-bond insert on $U$.  The interior-bond insert on $R$ is
    \begin{align}
        I^{\mathrm{int}}_{A,R,e,X}
         & =
        \frac{P_A(R)}{P_A(U)}
        \Ext_{U\subseteq R}(I_{A,U,e,X}).
        \notag
    \end{align}
    Thus the genuine tensors at the vertices of $R\setminus U$ are adjoined
    to a boundary insertion of $X$ on $e$, with the normalization changed
    from the non-boundary bond product of $U$ to that of $R$.
\end{definition}

\begin{theorem}[Closed state of an interior-bond insert]%
    \label{thm:peps_interior_bond_inserted_state}
    \lean{TNLean.PEPS.deformedRegionStateAssembled_interiorBondInserted_eq_smul_edgeInsertedCoeff}
    \leanok
    \uses{def:peps_interior_bond_inserted_region_insert, thm:peps_corner_extension_state_preserves, thm:peps_deformedRegionStateAssembled_bondInserted_eq_smul_edgeInsertedCoeff}
    Suppose all bond dimensions are positive.  For every global physical
    configuration $\omega$,
    \begin{align}
        \Psi^{\mathrm{as}}_{A,R,I^{\mathrm{int}}_{A,R,e,X}}(\omega)
         & =P_A(R)E_A(e,\omega,X).
        \notag
    \end{align}
\end{theorem}
\begin{proof}
    \leanok
    Corner extension preserves the closed state.  On the singleton $U$, the
    boundary orientation is the identity because $U$ contains the ordered
    left endpoint of $e$.  Hence its boundary insert has assembled state
    $P_A(U)E_A(e,\omega,X)$.  The factor $P_A(R)/P_A(U)$ cancels $P_A(U)$
    and gives the displayed identity.
\end{proof}

\begin{definition}[Staircase family for one virtual operation]%
    \label{def:peps_staircase_virtual_operation_insert}
    \lean{TNLean.PEPS.staircaseVirtualOperationInsert}
    \leanok
    \uses{def:peps_bond_inserted_region_insert, def:peps_interior_bond_inserted_region_insert, lem:peps_reference_edge_boundary_windows, lem:peps_reference_edge_interior_windows}
    Let $e$ be the horizontal reference edge and let
    $X\in\MN{D_B(e)}$.  Define an insert $C_j(X)$ on each staircase window
    $W_j$, for $0\le j<L+K$, by
    \begin{align}
        C_j(X)
         & =
        \begin{cases}
            I_{B,W_0,e,X^{\mathsf T}},    & j=0,        \\
            I^{\mathrm{int}}_{B,W_j,e,X}, & 1\le j<L,   \\
            I_{B,W_j,e,X},                & L\le j<L+K.
        \end{cases}
        \notag
    \end{align}
    The first window contains only the ordered right endpoint of $e$, the
    middle windows contain both endpoints, and the remaining windows contain
    only the ordered left endpoint.  The transpose in the first case makes
    all three boundary orientations represent the same matrix $X$.
\end{definition}

\begin{theorem}[Physical operators on the staircase windows]%
    \label{thm:peps_staircase_virtual_operation_physical_realization}
    \lean{TNLean.PEPS.staircaseVirtualOperationPhysicalOp}
    \lean{TNLean.PEPS.staircaseVirtualOperationPhysicalOp_realizes}
    \leanok
    \uses{def:peps_staircase_virtual_operation_insert, def:peps_physical_operator_of_region_insert, thm:peps_physical_operator_of_region_insert_realizes}
    Suppose every $L\times K$ window is injective.  For every
    $0\leq j<L+K$, there is a physical operator $O_j(X)$ on the physical
    space of $W_j$ such that
    \begin{align}
        O_j(X)\bigl(T_{B,W_j}(\mu)\bigr) & =C_j(X)(\mu)
        \qquad (\mu\in\partial W_j).
        \notag
    \end{align}
    Hence the virtual operation is realized locally on each window, before
    any contraction with the complementary tensor network.
\end{theorem}
\begin{proof}
    \leanok
    Each staircase window is an $L\times K$ translate and is therefore
    injective.  Apply the open-boundary realization theorem to $C_j(X)$.
\end{proof}

\begin{lemma}[The two end inserts of the staircase family]%
    \label{lem:peps_staircase_virtual_operation_end_inserts}
    \lean{TNLean.PEPS.staircaseVirtualOperationInsert_zero}
    \lean{TNLean.PEPS.staircaseVirtualOperationInsert_last}
    \leanok
    \uses{def:peps_staircase_virtual_operation_insert}
    The end members of the family are
    \begin{align}
        C_0(X)       & =I_{B,W_0,e,X^{\mathsf T}},
                     &
        C_{L+K-1}(X) & =I_{B,W_{L+K-1},e,X}.
        \notag
    \end{align}
\end{lemma}
\begin{proof}
    \leanok
    The first identity is the $j=0$ case of the definition.  Since $K>0$,
    the last index satisfies $L\le L+K-1$, which gives the second identity.
\end{proof}

\begin{theorem}[Common state of the staircase virtual-operation family]%
    \label{thm:peps_staircase_virtual_operation_common_state}
    \lean{TNLean.PEPS.deformedRegionStateAssembled_staircaseVirtualOperationInsert}
    \leanok
    \uses{def:peps_staircase_virtual_operation_insert, lem:peps_staircase_virtual_operation_end_inserts, lem:peps_reference_edge_left_endpoint_staircase_window, thm:peps_interior_bond_inserted_state, thm:peps_deformedRegionStateAssembled_bondInserted_eq_smul_edgeInsertedCoeff, thm:peps_regionInteriorBondProd_staircaseWindow_eq}
    Let $B$ be translation invariant and suppose all its bond dimensions are
    positive.  Then for $0\le j<L+K$ and every global physical configuration
    $\omega$,
    \begin{align}
        \Psi^{\mathrm{as}}_{B,W_j,C_j(X)}(\omega)
         & =P_B(W_0)E_B(e,\omega,X).
        \notag
    \end{align}
    In particular, all $L+K$ inserts have one common deformed state.
\end{theorem}
\begin{proof}
    \leanok
    On $W_0$, the boundary orientation sends $X^{\mathsf T}$ to $X$.
    On $W_1,\ldots,W_{L-1}$, apply the interior-bond identity.  On the
    remaining windows the ordered left endpoint lies in the window, so the
    boundary orientation fixes $X$.  Each case gives
    $P_B(W_j)E_B(e,\omega,X)$.  Translation invariance identifies
    $P_B(W_j)$ with $P_B(W_0)$.
\end{proof}

\begin{theorem}[Common staircase state transferred between tensors]%
    \label{thm:peps_staircase_physical_operation_same_state}
    \lean{TNLean.PEPS.deformedRegionStateAssembled_staircasePhysicalOp_sameState}
    \leanok
    \uses{thm:peps_physical_operator_same_state_transfer,
        thm:peps_staircase_virtual_operation_physical_realization,
        thm:peps_staircase_virtual_operation_common_state,
        thm:peps_regionInteriorBondProd_staircaseWindow_eq}
    Let $L,K>0$, and choose staircase coordinates $(a,b)$ with
    $a+2L\leq\mathrm{width}$ and $2K\leq\mathrm{height}$. Let $A$ and $B$
    be translation invariant tensors with the same physical dimension and the
    same closed state. Suppose every cyclic $L\times K$ window of $A$ is
    injective and $D_A(g)>0$ for every edge $g$. For
    $X\in\MN{D_A(e)}$, let $O_j^A(X)$ be the physical operators obtained from
    the staircase insert family of $A$, and define inserts on the windows of
    $B$ by
    \begin{align}
        C_j^{A\to B}(X)(\mu)
         & =O_j^A(X)\bigl(T_{B,W_j}(\mu)\bigr).
        \notag
    \end{align}
    Then, for all $j,j'\in\mathbb{N}$,
    \begin{align}
        \Psi^{\mathrm{as}}_{B,W_j,C_j^{A\to B}(X)}
         & =\Psi^{\mathrm{as}}_{B,W_{j'},C_{j'}^{A\to B}(X)}.
        \notag
    \end{align}
    The source family has $0\leq j<L+K$. In the definition used here the
    truncated offsets stabilize at $W_{L+K-1}$ for $j\geq L+K-1$; the indices
    are not read modulo $L+K$. No equality between the edge-dependent bond
    dimensions of $A$ and $B$, and no positivity assumption on those of $B$,
    is required.
\end{theorem}
\begin{proof}
    \leanok
    For each $j$, (\ref{eq:peps_physical_op_same_state}) and the
    open-boundary realization of $O_j^A(X)$ give
    \begin{align}
        P_B(W_j)\Psi^{\mathrm{as}}_{A,W_j,C_j^A(X)}
         & =P_A(W_j)\Psi^{\mathrm{as}}_{B,W_j,C_j^{A\to B}(X)}.
        \notag
    \end{align}
    The staircase family of $A$ has one common deformed state. Translation
    invariance gives $P_A(W_j)=P_A(W_0)$ and
    $P_B(W_j)=P_B(W_0)$. Hence the right-hand side is independent of $j$.
    Positivity of the bond dimensions makes $P_A(W_0)$ nonzero, so it may be
    cancelled.
\end{proof}

\begin{definition}[The left staircase end operation $O_1^A(X)$]%
    \label{def:peps_staircase_O1_physical_operation}
    \lean{TNLean.PEPS.staircaseO1PhysicalOp}
    \leanok
    \uses{thm:peps_staircase_virtual_operation_physical_realization}
    For $X\in\MN{D_A(e)}$, define
    $O_1^A(X)=O_{L+K-1}^A(X)$ after identifying the last staircase
    window with the left end window.
\end{definition}

\begin{definition}[The underlying right staircase end operation $O_3^A(X)$]%
    \label{def:peps_staircase_O3_physical_operation}
    \lean{TNLean.PEPS.staircaseO3PhysicalOp}
    \leanok
    \uses{thm:peps_staircase_virtual_operation_physical_realization}
    For $X\in\MN{D_A(e)}$, define $O_3^A(X)=O_0^A(X)$ after identifying
    the first staircase window with the right end window.
    The relation (\ref{eq:peps_staircase_o3_cross_tensor_relation}) has the
    paper's transposed $O_3^{\mathsf T}$ orientation; no transpose is chosen
    on the abstract physical endomorphism space.
\end{definition}

\begin{theorem}[Cross-tensor virtual operation on the staircase edge]%
    \label{thm:peps_cross_tensor_virtual_operation_of_staircase}
    \lean{TNLean.PEPS.existsUnique_crossTensorVirtualOperation_of_staircasePhysicalOp_sameState}
    \leanok
    \uses{def:peps_staircase_O1_physical_operation,
        def:peps_staircase_O3_physical_operation}
    Let $L,K\geq2$, and choose the displayed non-wrapping horizontal
    staircase coordinates $(a,b)$ with
    \begin{align}
        1    & \leq a,
             & a+2L                 & \leq \mathrm{width},
             & b+2K-1               & \leq \mathrm{height},
        \notag                                              \\
        2L+1 & \leq \mathrm{width},
             & 2K+1                 & \leq \mathrm{height}.
        \notag
    \end{align}
    Suppose that $A$ and $B$ are translation invariant, generate the same
    closed state, every cyclic $L\times K$ window of either tensor is
    injective, and every bond dimension of either tensor is positive. For
    $X\in\MN{D_A(e)}$, there is a unique $Y\in\MN{D_B(e)}$ such that, for
    every external boundary configuration $\eta_{\mathrm L}$ and
    $\eta_{\mathrm R}$, and every reference-edge index $k$,
    \begin{align}
        O_1^A(X)\bigl(T_{B,W_{\mathrm L}}(k,\eta_{\mathrm L})\bigr)
         & =\sum_jY_{jk}T_{B,W_{\mathrm L}}(j,\eta_{\mathrm L}),
        \label{eq:peps_staircase_o1_cross_tensor_relation}       \\
        O_3^A(X)\bigl(T_{B,W_{\mathrm R}}(k,\eta_{\mathrm R})\bigr)
         & =\sum_jY_{kj}T_{B,W_{\mathrm R}}(j,\eta_{\mathrm R}).
        \label{eq:peps_staircase_o3_cross_tensor_relation}
    \end{align}
    Thus the relation (\ref{eq:peps_staircase_o3_cross_tensor_relation}) has the
    paper's $O_3^{\mathsf T}$ orientation.
    No equality between $D_A(e)$ and $D_B(e)$ is assumed or concluded, and
    this statement contains no multiplication or gauge assertion.
\end{theorem}
\begin{proof}
    \leanok
    \uses{thm:peps_staircase_physical_operation_same_state,
        thm:peps_staircase_pair_insert_eq_open,
        thm:peps_existsUnique_virtualOperation_of_horizontal_staircase_end_pair,
        thm:peps_regionInjectivityUnionClosure_of_overlap}
    The transferred physical inserts form one common $B$-state. The
    overlapping-union lemma supplies the union closure needed to compare
    consecutive staircase windows, and corner completion gives equality of
    the two end inserts with open boundary. At the two ends these inserts are
    precisely the actions of $O_3^A(X)$ and $O_1^A(X)$ on the genuine
    $B$-window tensors. Injectivity of those two windows and the two-end
    comparison give the displayed relations and the uniqueness of $Y$.
\end{proof}

\begin{theorem}[Composition of the canonical staircase end operations]%
    \label{thm:peps_staircase_end_physical_operation_multiplication}
    \lean{TNLean.PEPS.staircaseO1PhysicalOp_mul}
    \lean{TNLean.PEPS.staircaseO3PhysicalOp_mul}
    \leanok
    \uses{def:peps_staircase_O1_physical_operation,
        def:peps_staircase_O3_physical_operation}
    Let $L,K>0$, and choose staircase coordinates $(a,b)$ with
    $a+2L\leq\mathrm{width}$ and $2K\leq\mathrm{height}$. Suppose every
    cyclic $L\times K$ window of $A$ is injective. For matrices
    $X,Y\in\MN{D_A(e)}$, the canonical end operations satisfy
    \begin{align}
        O_1^A(XY) & =O_1^A(X)O_1^A(Y),
        \notag                         \\
        O_3^A(XY) & =O_3^A(Y)O_3^A(X).
        \notag
    \end{align}
    Thus the right relation preserves multiplication in ordinary order after
    passing to the orientation $(O_3^A)^{\mathsf T}$ used in Lemma~5.
\end{theorem}
\begin{proof}
    \leanok
    \uses{lem:peps_staircase_virtual_operation_end_inserts,
        thm:peps_lemma5_end_operation_multiplication}
    Split the virtual boundary of either end window into the reference-edge
    coordinate and the remaining external coordinates.  Let $E$ be this
    splitting map, and write $C_M$ for insertion of $M$ on the reference edge.
    On every fixed external fibre,
    \begin{align}
        C_M\bigl(E^{-1}(j,\eta)\bigr)
         & =\sum_i M_{ij}\,T_{A,W}\bigl(E^{-1}(i,\eta)\bigr).
        \notag
    \end{align}
    The chosen left inverse of the injective blocked tensor is a left inverse
    on the whole boundary-coordinate space, so this identity gives
    $O_W(M)O_W(N)=O_W(MN)$.  The first window carries $X^{\mathsf T}$; hence
    $(XY)^{\mathsf T}=Y^{\mathsf T}X^{\mathsf T}$ reverses the order of the
    underlying right operations.
\end{proof}

\begin{definition}[Canonical cross-tensor matrix assignment]%
    \label{def:peps_staircase_cross_tensor_virtual_operation}
    \lean{TNLean.PEPS.staircaseCrossTensorVirtualOperation}
    \lean{TNLean.PEPS.staircaseCrossTensorVirtualOperation_spec}
    \leanok
    \uses{thm:peps_cross_tensor_virtual_operation_of_staircase}
    Under the hypotheses of
    Theorem~\ref{thm:peps_cross_tensor_virtual_operation_of_staircase}, in the
    paper's notation the input and uniquely recovered matrices are
    called $X$ and $Y$.  Writing $Y_{A\to B}(X)$ only to distinguish the two
    tensor families, define
    \begin{align}
        Y_{A\to B}(X) & \in\MN{D_B(e)}
        \notag
    \end{align}
    to be the unique matrix satisfying
    (\ref{eq:peps_staircase_o1_cross_tensor_relation}) and
    (\ref{eq:peps_staircase_o3_cross_tensor_relation}) for every left and
    right external boundary configuration.  This definition makes no choice
    of an external boundary configuration.
\end{definition}

\begin{theorem}[Multiplicativity of the cross-tensor assignment]%
    \label{thm:peps_staircase_cross_tensor_virtual_operation_multiplication}
    \lean{TNLean.PEPS.staircaseCrossTensorVirtualOperation_mul}
    \leanok
    \uses{def:peps_staircase_cross_tensor_virtual_operation}
    Under the hypotheses of
    Theorem~\ref{thm:peps_cross_tensor_virtual_operation_of_staircase}, for all
    $X_1,X_2\in\MN{D_A(e)}$,
    \begin{align}
        Y_{A\to B}(X_1X_2)
         & =Y_{A\to B}(X_1)Y_{A\to B}(X_2).
        \notag
    \end{align}
\end{theorem}
\begin{proof}
    \leanok
    \uses{thm:peps_staircase_end_physical_operation_multiplication,
        thm:peps_lemma5_end_operation_multiplication}
    Set $Y_r=Y_{A\to B}(X_r)$ for $r=1,2$.  For every left external
    boundary configuration, the two $O_1$ relations compose in ordinary order:
    \begin{align}
        O_1^A(X_1)\Bigl(O_1^A(X_2)
        \bigl(T_{B,W_{\mathrm L}}(k,\eta_{\mathrm L})\bigr)\Bigr)
         & =\sum_i (Y_1Y_2)_{ik}\,
        T_{B,W_{\mathrm L}}(i,\eta_{\mathrm L}).
        \notag
    \end{align}
    For every right external boundary configuration, the underlying $O_3$
    operations compose in reverse order:
    \begin{align}
        O_3^A(X_2)\Bigl(O_3^A(X_1)
        \bigl(T_{B,W_{\mathrm R}}(k,\eta_{\mathrm R})\bigr)\Bigr)
         & =\sum_j (Y_1Y_2)_{kj}\,
        T_{B,W_{\mathrm R}}(j,\eta_{\mathrm R}).
        \notag
    \end{align}
    This is ordinary multiplication in the $O_3^{\mathsf T}$ orientation.  The
    preceding theorem identifies the left sides with the canonical operations
    for $X_1X_2$.  Both end relations therefore realize $Y_1Y_2$, and global
    uniqueness identifies this product with $Y_{A\to B}(X_1X_2)$.
\end{proof}

\begin{theorem}[End-window identity for one virtual operation]%
    \label{thm:peps_end_window_identity_staircase_virtual_operation}
    \lean{TNLean.PEPS.regionInsertedCoeff_endWindows_eq_staircaseVirtualOperation}
    \leanok
    \uses{lem:peps_assembleRegionSigma_restrict, lem:peps_reference_edge_left_endpoint_staircase_window, thm:peps_regionInsertedCoeff_eq_smul_edgeInsertedCoeff, thm:peps_regionInteriorBondProd_staircaseWindow_eq}
    Suppose the tensor is translation invariant and the staircase satisfies
    the stated size bounds.  Then for every $X\in\MN{D_B(e)}$ and every
    global physical configuration $\omega$,
    \begin{align}
        C_B(W_0,e,X^{\mathsf T};\omega|_{W_0},
        \omega|_{V\setminus W_0})
         & =
        C_B(W_{L+K-1},e,X;\omega|_{W_{L+K-1}},
        \omega|_{V\setminus W_{L+K-1}}).
        \notag
    \end{align}
\end{theorem}
\begin{proof}
    \leanok
    The region-to-edge identity sends both sides to the edge coefficient
    with $X$ on the same reference edge: the first boundary orientation
    transposes $X^{\mathsf T}$, while the last fixes $X$.  Translation
    invariance identifies the two non-boundary bond products.
\end{proof}
